\section{Fuente de Datos y Consideraciones Éticas}

\subsection{Origen de los Datos y Unidad de Observación}
Los datos utilizados en este estudio provienen de la \textit{National Health and Nutrition Examination Survey} (NHANES), un programa de estudios diseñado para evaluar el estado de salud y nutrición de adultos y niños en los Estados Unidos. Específicamente, se utiliza el ciclo de datos continuos correspondiente al período 2021-2023. La unidad de observación es una persona individual examinada (NCHS, 2024). 

La base de datos original se distribuye en múltiples módulos independientes en formato SAS Transport (\texttt{.xpt}). Para la consolidación de la muestra analítica, se cruzaron los registros de cuatro módulos clave utilizando el identificador único de cada paciente (\texttt{SEQN}): Demografía (\texttt{DEMO\_L}), Medidas Corporales (\texttt{BMX\_L}), Presión Arterial (\texttt{BPXO\_L}) y Colesterol Total (\texttt{TCHOL\_L}).

\subsection{Consideraciones Éticas}
El uso de esta base de datos se enmarca en un contexto estrictamente académico y estadístico. Los datos de NHANES son de dominio público y han sido previamente anonimizados por el \textit{National Center for Health Statistics} (NCHS) del CDC, garantizando que no sea posible identificar a los individuos encuestados. El análisis propuesto no busca la estigmatización de ningún subgrupo demográfico, sino comprender las relaciones entre indicadores antropométricos y cardiovasculares.
